% Public website CV derived from the application master.
\documentclass[letterpaper,11pt]{article}

\usepackage{latexsym}
\usepackage[empty]{fullpage}
\usepackage{titlesec}
\usepackage{marvosym}
\usepackage[usenames,dvipsnames]{color}
\usepackage{verbatim}
\usepackage{enumitem}
\usepackage[pdftex]{hyperref}
\usepackage{fancyhdr}

\usepackage{etaremune}
\usepackage{enumitem}
\usepackage{xcolor}
\hypersetup{
    colorlinks=true,
    linkcolor=MidnightBlue,
    urlcolor=MidnightBlue,
    citecolor=MidnightBlue
}

\pagestyle{fancy}
\fancyhf{} % clear all header and footer fields
\fancyfoot{}
\renewcommand{\headrulewidth}{0pt}
\renewcommand{\footrulewidth}{0pt}

% Adjust margins
\addtolength{\oddsidemargin}{-0.375in}
\addtolength{\evensidemargin}{-0.375in}
\addtolength{\textwidth}{1in}
\addtolength{\topmargin}{-.5in}
\addtolength{\textheight}{1.0in}

\urlstyle{same}

\raggedbottom
\raggedright
\setlength{\tabcolsep}{0in}

% Sections formatting
\titleformat{\section}{
  \vspace{-4pt}\scshape\raggedright\large
}{}{0em}{}[\color{black}\titlerule \vspace{-5pt}]

%-------------------------
% Custom commands
\newcommand{\resumeItem}[2]{
  \item\small{
    \textbf{#1}{: #2 \vspace{-2pt}}
  }
}

\newcommand{\resumeSubheading}[4]{
  \vspace{-1pt}\item
    \begin{tabular*}{0.97\textwidth}{l@{\extracolsep{\fill}}r}
      \textbf{#1} & #2 \\
      \textit{\small#3} & \textit{\small #4} \\
    \end{tabular*}\vspace{-5pt}
}

\newcommand{\resumeSubItem}[2]{\resumeItem{#1}{#2}\vspace{-4pt}}

\renewcommand{\labelitemii}{$\circ$}

\newcommand{\resumeSubHeadingListStart}{\begin{itemize}[leftmargin=*]}
\newcommand{\resumeSubHeadingListEnd}{\end{itemize}}
\newcommand{\resumeItemListStart}{\begin{itemize}}
\newcommand{\resumeItemListEnd}{\end{itemize}\vspace{-5pt}}

%-------------------------------------------
%%%%%%  CV STARTS HERE  %%%%%%%%%%%%%%%%%%%%%%%%%%%%


\begin{document}

%----------HEADING-----------------
\begin{tabular*}{\textwidth}{l@{\extracolsep{\fill}}r}
  \textbf{\Large Yimeng Tang} & Email: \href{mailto:ymtang@ucsc.edu}{ymtang@ucsc.edu}\\
  {University of California, Santa Cruz} & Website: \href{https://ymtang2000.github.io/}{ymtang2000.github.io} \\
  {Department of Astronomy \& Astrophysics} & ORCID: \href{https://orcid.org/0000-0003-2876-577X}{0000-0003-2876-577X}\\
  {1156 High St, Santa Cruz, CA 95064} & {}\\
\end{tabular*}

\section{Research Interests}
\begin{enumerate}[itemsep=-3pt]
  %\item[-] Assembly histories of galaxies and their star clusters across diverse environments
  %\item[-] Ultra-diffuse galaxies and the low-surface-brightness Universe
  %\item[-] Dark matter deficient galaxies, galaxy--halo connection
  %\item[-] Stellar feedback and starburst galaxies
  %\item[-] Co-evolution of galaxies and supermassive black holes
  \item[-] Formation and evolution of dwarf galaxies and low-surface-brightness galaxies
  \item[-] Globular cluster systems as tracers of galaxy assembly
  \item[-] Stellar populations and star formation histories of galaxies and star clusters
  \item[-] Stellar feedback, gas flows, and environmental transformation of galaxies
  \item[-] Dark-matter-deficient galaxies and galaxy–halo connection
  \item[-] Emission-line diagnostics of star formation and nuclear activity
\end{enumerate}

%-----------EDUCATION-----------------
\section{Education}
  {{\bf Ph.D. in Astronomy}, University of California, Santa Cruz} \hfill {Jun 2027 (expected)}\\
  \vspace{0.05cm}
  \hspace{0.3cm} {\em Advisors: Aaron Romanowsky, Kevin Bundy}\\
  \hspace{0.3cm} {\em Thesis:  Decoding galaxy formation through the diversity of dwarf galaxies and their globular clusters}\\
  \vspace{0.2cm}
  {{\bf M.S. in Astronomy}, University of California, Santa Cruz} \hfill {2024}\\
  \vspace{0.2cm}
  {{\bf B.S. in Astronomy (Highest Honor)}, University of Science and Technology of China} \hfill {2021}\\
  \vspace{0.05cm}
  \hspace{0.3cm} {\em Advisors: Hong-Xin Zhang, Xu Kong}\\
  \hspace{0.3cm} {\em Thesis: Observational constraints on the evolution of nearby galaxies from IFU spectroscopy}\\

\section{Selected Awards \& Honors}
  Chancellor's Dissertation Year Fellowship, {\small UCSC} \hfill{2026}\\
  USTC Undergraduate Honorary Rank {\small($\sim$ top 5\%)}, {\small USTC} \hfill{2021}\\
  \vspace{0.03cm}Guo Moruo Scholarship {\small($\sim$ top 2\%, {\bf highest honor for undergraduates at USTC})}, {\small USTC} \hfill{2020}\vspace{0.03cm}\\ 
  \vspace{0.03cm}National Scholarship {\small($\sim$ top 2\%)}, {\small Ministry of Education of China} \hfill{2020~{\footnotesize\&}~2018}\\ 
  \vspace{0.03cm}NAOC Scholarship, {\small National Astronomical Observatories of Chinese Academy of Sciences} \hfill{2020}\\
  \vspace{0.03cm}Scholarship of Yan Jici Talent Program in Physics {\small($\sim$ top 10\%)}, {\small USTC} \hfill{2020~{\footnotesize\&}~2019}\\ 
  \vspace{0.03cm}Outstanding Student Scholarship, Gold Prize {\small($\sim$ top 5\%)}, {\small USTC} \hfill{2019}\\ 
  \vspace{0.03cm}Outstanding Freshman Scholarship, {\small USTC} \hfill{2017}

\section{Observing programs}
{\bf As Principal Investigator:}
\vspace{-0.2cm}
\begin{etaremune}[itemsep=0pt]
\item Keck 2025A (0.5 night, KCWI): {\em Spectroscopic mapping of an ultra-diffuse galaxy with massive star-forming clumps}
\end{etaremune}
\vspace{-0.2cm}
{\bf As Co-Investigator with Major Contributions:}
\vspace{-0.2cm}
\begin{etaremune}[itemsep=0pt]
\item Roman Cycle 1 (Analysis): {\em Extragalactic globular clusters and galaxy formation with Roman}
\item Subaru 2026B (6 hours, HSC): {\em Imaging of Perseus cluster dwarf galaxies and globular clusters}
\item JWST Cycle 5 (10.3 hours, NIRCam \& NIRISS): {\em Testing the formation of a trail of dwarf galaxies through tip of the red giant branch distances}
\item Keck 2026A (1.5 nights, KCWI): {\em The stellar population gradients of ultra-diffuse galaxies}
\item VLT 2025A (3 hours, MUSE): {\em Young and Cluster-Rich: The Progenitors of Globular Cluster-Rich Ultra-Diffuse Galaxies?}
\item Keck 2025A (2 nights, KCWI): {\em The stellar population gradients of ultra-diffuse galaxies}
\item VLT 2024B (17 hours, MUSE): {\em FCC 224: another dark matter free dwarf galaxy?} 
\item Keck 2024B (1 night, KCWI): {\em Dark matter in NGC 1052 dwarf galaxies}
\item Keck 2024B (2 nights, KCWI): {\em Globular clusters and dark matter in ultra-diffuse galaxies}
\item Keck 2024B (1 night, KCWI): {\em A critical test for dark matter in dwarf galaxies of the NGC 1052 group}
\item Keck 2023A (1 night, DEIMOS): {\em Unravelling the origins of cluster ultra-diffuse galaxies}
\item Keck 2022B (3 nights, DEIMOS): {\em Unravelling the origins of ultra-diffuse galaxies in the Perseus cluster}
\end{etaremune}
\vspace{-0.2cm}
{\bf Additional Co-Investigator Programs:}
\vspace{-0.2cm}
\begin{etaremune}[itemsep=0pt]
\item Gemini 2026B (13 hours, GMOS): {\em A Benchmark Globular Cluster System for the Survey Era: GMOS Spectroscopy of NGC 1527}
\item Keck 2026B (2.5 nights, KCWI \& LRIS): {\em Testing a Possible Second Trail of Dark-Matter-Deficient Dwarf Galaxies in the Outskirts of Fornax}
\item HST Cycle 29 (15 orbits, ACS): {\em A trail of dark matter-free galaxies in the NGC 1052 group}
\end{etaremune}
{\bf Observing Experience: }{Keck Observatory (12 nights, KCWI, DEIMOS, OSIRIS), Lick Observatory (3 nights, Shane, Nickel)}

\section{Selected Talks \& Presentations}
\begin{enumerate}[itemsep=0pt]
  \item[-] Dept. of Astronomy, Harvard University {\small (Cambridge, Massachusetts, USA)} \hfill {Sep 2026}\\
  {\em Talk: Stellar population modeling of dwarf galaxies and globular clusters with Rubin, Roman and Euclid}
  \item[-] Dept. of Astronomy, Yale University {\small (New Haven, Connecticut, USA)} \hfill {Sep 2026}\\
  {\em Talk: Puzzles of the Small: Dwarf Galaxies and Their Globular Clusters Across Environments}
  \item[-] 2026 Santa Cruz Galaxy Workshop {\small (Santa Cruz, California, USA)} \hfill {Aug 2026}\\
  {\em Talk: Stellar Populations of Dwarf Galaxies in the Local Universe}\\
  {\em Talk: Extreme Dwarf Galaxies in the Local Universe}\\
  \item[-] 2026 Rubin Observatory Community Workshop {\small (SLAC, Menlo Park, California, USA)} \hfill {Jul 2026}\\
  {\em Talk: A Comprehensive Study of Dwarf Galaxies in the Era of Next-generation Telescopes}\\
  {\em Talk: Globular Clusters in Nearby Galaxies: Numbers and Stellar Populations}\\
  \item[-] Near Field Cosmology in the Era of Big Data: Local Group and Beyond {\small (Toronto, Canada)} \hfill {Jul 2026}\\
  {\em Talk: A Comprehensive Study of Dwarf Galaxies with Euclid, Roman and Rubin}\\
  \item[-] Kavli Institute for Astronomy and Astrophysics, Peking University {\small (Beijing, China)} \hfill {Apr 2026}\\
  Dept. of Astronomy, Tsinghua University {\small (Beijing, China)}\\
  Dept. of Astronomy, University of Science and Technology of China {\small (Hefei, China)}\\
  School of Astronomy \& Space Science, Nanjing University {\small (Nanjing, China)}\\
  {\em Talk: Puzzles of the Small: Diversity of Dwarf Galaxies and Their Globular Clusters Across Environments}
  \item[-] 10 Years of Ultra-Diffuse Galaxies: Past Perspectives and Future Prospects {\small (Benasque, Spain)} \hfill {Aug 2025}\\
  {\em Talk: Globular Clusters as Guides: UDGs and Dwarf Galaxies in Various Environments}
  \item[-] 2025 Santa Cruz Galaxy Workshop {\small (Santa Cruz, California, USA)} \hfill {Aug 2025}\\
  {\em Talk: Connection Between Dwarf Galaxies and Globular Clusters -- Insights from the Perseus Cluster}\\
  {\em Talk: Central Emission-line Spectra of SDSS Galaxies: The Basic Evolutionary Path and its Relation to AGN Activity}
  \item[-] 2024 Santa Cruz Galaxy Workshop {\small (Santa Cruz, California, USA)} \hfill {Aug 2024}\\
  {\em Talk: A New Candidate Dark Matter Deficient Galaxy with Unusual Globular Clusters}
  \item[-] Dwarf Galaxies, Star Clusters, and Streams in the LSST Era {\small (Chicago, Illinois, USA)} \hfill {Jul 2024}\\
  {\em Poster: Observational Constraints on the Formation of Dark-matter Deficient Galaxies from Morphologies, Stellar Populations and Globular Clusters}
  \item[-] Dept. of Astronomy, University of Science and Technology of China {\small (Hefei, China)} \hfill {Dec 2023}\\
  {\em Talk: Testing the Bullet Dwarf Collision Scenario in the NGC 1052 Group}
  \item[-] 2023 Santa Cruz Galaxy Workshop {\small (Santa Cruz, California, USA)} \hfill {Aug 2023}\\
  {\em Talk: Testing the Bullet Dwarf Collision Scenario in the NGC 1052 Group}
  \item[-] UC Santa Cruz FLASH Talks {\small (Santa Cruz, California, USA)} \hfill {Jun 2023}\\
  {\em Talk: Testing the Bullet Dwarf Collision Scenario in the NGC 1052 Group}
  \item[-] Yan Jici Physics Symposium, University of Science and Technology of China {\small (Hefei, China)} \hfill {May 2021}\\
  {\em Talk: Observational constraints on the evolution of nearby galaxies from IFU spectroscopy}
  \item[-] 2019 Annual Conference of the Chinese Astronomical Society {\small (Delingha, China)} \hfill {Sep 2019}\\
  {\em Poster: The Origin of Surface Brightness Profile Breaks of Disk Galaxies from MaNGA}
\end{enumerate}

\section{Press coverage}
\begin{enumerate}[itemsep=0pt]
  \item[-] Scientific American: \href{https://www.scientificamerican.com/article/galaxies-without-dark-matter-mystify-astronomers/}{Galaxies without dark matter mystify astronomers}
  \item[-] Keck Observatory: \href{https://keckobservatory.org/dmf-dwarf-galaxy/}{Evidence of Dark Matter-Free Dwarf Galaxy Challenges Conventional Galaxy Formation Models}
  \item[-] Phys.org: \href{https://phys.org/news/2022-10-star-formation-history-galaxy-ngc.html}{Observations shed more light on the star formation history of galaxy NGC 2915}
  %\item[-] Space: \href{https://www.space.com/the-universe/ghostly-galaxy-without-dark-matter-baffles-astronomers}{Ghostly galaxy without dark matter baffles astronomers}
\end{enumerate}

\section{Service, Membership \& Teaching}
Member, Vera C. Rubin Observatory LSST Dark Energy Science Collaboration\\
Member, Roman Science Collaboration\\
Member, Analysis of Galaxies At The Extremes (AGATE) Collaboration\\
Mentor, Introduction to Research in Physics and Astrophysics (ASTR 9, UCSC) \hfill {2026 Winter~\&~Spring}\\
Teaching Assistant, Introduction to Scientific Computing (ASTR 119, UCSC) \hfill {2026 Winter}\\
Teaching Assistant, Practical Programming for the Sciences (ASTR 19, UCSC) \hfill {2024~\&~2022 Spring}\\
Graduate Admissions Committee Member, UCSC \hfill {2023-2024}\\
Prospective Student Visiting Committee Member, UCSC \hfill {2022~\&~2023}\\
Student Admissions Ambassador, USTC \hfill {2021}\\
Teaching Assistant, Astronomical Labs (ASTR 167, USTC) \hfill {2020 Fall}\\
Teaching Assistant, Observational Astronomy (ASTR 170, USTC) \hfill {2020 Spring}


\section{References}
{\bf Aaron Romanowsky} (UC Santa Cruz \& San Jos\'e State U) \hfill {aaron.romanowsky@sjsu.edu}\\
\vspace{0.1cm}
{\bf Pieter van Dokkum} (Yale) \hfill {pieter.vandokkum@yale.edu}\\
\vspace{0.1cm}
{\bf Kevin Bundy} (UC Santa Cruz) \hfill {kbundy@ucsc.edu}\\
\vspace{0.1cm}
{\bf Sandra Faber} (UC Santa Cruz) \hfill {faber@ucolick.org}\\
\vspace{0.1cm}
{\bf Hong-Xin Zhang} (University of Science and Technology of China) \hfill {hzhang18@ustc.edu.cn}\\
\vspace{0.1cm}
{\bf Xu Kong} (University of Science and Technology of China) \hfill {xkong@ustc.edu.cn}

\section{Publications}
\href{https://ui.adsabs.harvard.edu/public-libraries/wkkyIxzqSveeiQVd0Qn-HQ}{Link to my ADS publication library}\\
\vspace{0.1cm}
{\bf First-author Publications:}
\vspace{-0.15cm}
\begin{etaremune}[itemsep=-2pt]
\item {\bf Y. Tang}, A. J. Romanowsky, et al., {Toward Extragalactic Globular Cluster Stellar Populations: Optical+NIR SED Fitting and Initial Results in Dwarf Galaxies, 2026, to be submitted}
\item {\bf Y. Tang}, A. J. Romanowsky, N. Okabe, and K. A. Bundy, \href{https://arxiv.org/abs/2611.01653}{An Unusually Clumpy Ultra-Diffuse Galaxy in the Coma Cluster: A Missing Link for Globular Cluster Formation? 2026, submitted to ApJ (arXiv:2611.01653)}
\item {\bf Y. Tang}, Y. Luo, C. Zhou, \href{https://arxiv.org/abs/2611.02580}{The Stellar-Mass Dependence of Feedback-Driven Outflows in DESI DR1 Star-Forming Galaxies, 2026, submitted to ApJL (arXiv:2611.02580)}
\item {\bf Y. Tang}, G. S. Anand, A. J. Romanowsky, et al., \href{https://ui.adsabs.harvard.edu/abs/2026ApJ..1005L..16T/abstract}{New Measurements of Distances to Galaxies in the NGC 1052 Field with the Hubble and James Webb Space Telescopes: Testing the Bullet-Dwarf Origin of the Trail, 2026, ApJL, 1005, L16}
\item {\bf Y. Tang}, A. J. Romanowsky, S. Huang, et al., \href{https://ui.adsabs.harvard.edu/abs/2026ApJ...998..254T/abstract}{Connection Between Dwarf Galaxies and Globular Clusters: Insights from the Perseus Cluster Using Subaru Imaging and Keck Spectroscopy, 2026, ApJ, 998, 254}
\item {\bf Y. Tang}, A. J. Romanowsky, J. S. Gannon, et al., \href{https://ui.adsabs.harvard.edu/abs/2025ApJ...982....1T/abstract}{An Unexplained Origin for the Unusual Globular Cluster System in the Ultra-diffuse Galaxy FCC 224, 2025, ApJ, 982, 1}
\item {\bf Y. Tang}, A. J. Romanowsky, P. G. van Dokkum, et al., \href{https://ui.adsabs.harvard.edu/abs/2025ApJ...978...21T/abstract}{Testing the Bullet Dwarf Collision Scenario in the NGC 1052 Group Through Morphologies and Stellar Populations, 2025, ApJ, 978, 21}
\item {\bf Y. Tang}, B. Tao, H.-X. Zhang, et al., \href{https://ui.adsabs.harvard.edu/abs/2022A%26A...668A.179T/abstract}{Unveiling the Formation of NGC 2915 with MUSE: A Counter-rotating Stellar Disk Embedded in a Disordered Gaseous Environment, 2022, A\&A, 668, A179}
\item {\bf Y. Tang}, Q. Chen, H.-X. Zhang, et al., \href{https://ui.adsabs.harvard.edu/abs/2020ApJ...897...79T/abstract}{New Constraints on the Origin of Surface Brightness Profile Breaks of Disk Galaxies from MaNGA, 2020, ApJ, 897, 79}
\end{etaremune}
{\bf Contributing-Author Publications:}
\vspace{-0.15cm}
\begin{etaremune}[itemsep=-2pt]
%Aaron Romanowsky -- Coma
%Wenhe Lyu -- SBF
\item Z. A. Coustier, A. J. Romanowsky, {\bf Y. Tang}, et al., {Stellar Populations of Star Clusters in the Local Group Dwarf Galaxy NGC 6822 Using Keck Spectroscopy, 2026, to be submitted}
\item A. Levitskiy, L. Valenzuela, D. A. Forbes, et al. (including {\bf Y. Tang}), {Diffuse Dwarf Galaxies in Galaxy Clusters: II. Stellar Kinematics, 2026, submitted}
\item D. Vaz, A. Levitskiy, D. A. Forbes, et al. (including {\bf Y. Tang}), \href{https://arxiv.org/abs/2610.07467}{The Rogue One: A Dark Matter-Dominated Dwarf in the Trail of Dark Matter-Deficient Galaxies, 2026, Accepted by MNRAS}
\item A. Levitskiy, D. A. Forbes, J. S. Gannon, {\bf Y. Tang}, et al., \href{https://ui.adsabs.harvard.edu/abs/2026MNRAS.551g1568L/abstract}{Diffuse Dwarf Galaxies in Galaxy Clusters: I. Stellar Populations and Radial Gradients, 2026, MNRAS, 551, 1568}
\item M. L. Buzzo, P. G. van Dokkum, M. Hilker, et al. (including {\bf Y. Tang}), \href{https://ui.adsabs.harvard.edu/abs/2026arXiv260524099B/abstract}{A Twin Set of Dark Matter-Deficient Dwarfs: FCC 224 and FCC 240 as Analogues of NGC 1052-DF2 and DF4, 2026, ApJ, 1005, 15}
\item D. A. Forbes, B. van Heumen, {\bf Y. Tang}, \href{https://ui.adsabs.harvard.edu/abs/2026MNRAS.548ag765F/abstract}{Revisiting the distance and the globular cluster system of the remarkable galaxy UDG1 in the NGC 5846 group, 2026, MNRAS, 548, 765}
\item D. Martínez-Delgado, A. J. Romanowsky, {\bf Y. Tang}, et al., \href{https://ui.adsabs.harvard.edu/abs/2026A%26A...708A.172M/abstract}{A Pearl in the Shell: an ultra-compact dwarf within the tidal debris surrounding spiral galaxy NGC 7531, 2026, A\&A, 708, A172}
\item M. A. Keim, P. G. van Dokkum, Z. Shen, et al. (including {\bf Y. Tang}), \href{https://ui.adsabs.harvard.edu/abs/2025ApJ...988..165K/abstract}{Kinematic Confirmation of a Remarkable Linear Trail of Galaxies in the NGC 1052 Field, Consistent with Formation in a High-Speed Bullet Dwarf Collision, 2025, ApJ, 988, 165}
\item  A. J. Romanowsky, {\bf Y. Tang}, K. A. Bundy, et al., \href{https://ui.adsabs.harvard.edu/abs/2025RNAAS...9..181R/abstract}{Morphology and Stellar Populations of a Candidate Ultra-diffuse Galaxy in Early Euclid and Rubin Imaging, 2025, RNAAS, 9, 181}
\item T. Li, H.-X. Zhang, W. Lyu, {\bf Y. Tang}, et al., \href{https://ui.adsabs.harvard.edu/abs/2025A%26A...698A.208L/abstract}{A Negative Stellar Mass-Gaseous Metallicity Gradient Relation of Dwarf Galaxies Modulated by Stellar Feedback, 2025, A\&A, 698, A208}
\item M. L. Buzzo, D. A. Forbes, A. J. Romanowsky, et al. (including {\bf Y. Tang}), \href{https://ui.adsabs.harvard.edu/abs/2025A%26A...695A.124B/abstract}{A New Class of Dark Matter-Free Dwarf Galaxies? I. Clues from FCC 224, NGC 1052-DF2 and NGC 1052-DF4,  2025, A\&A, 695, A124}
\item W. Lyu, H.-X. Zhang, S. Paudel, et al. (including {\bf Y. Tang}), \href{https://ui.adsabs.harvard.edu/abs/2025ApJ...979...85L/abstract}{Dependence of Metal Enrichment of Nuclear Star Clusters on Galaxy Stellar Mass, 2025, ApJ, 979, 85}
\item A. J. Romanowsky, S. S. Larsen, A. Villaume, et al. (including {\bf Y. Tang}), \href{https://ui.adsabs.harvard.edu/abs/2023MNRAS.518.3164R/abstract}{Low-density star cluster formation: discovery of a young faint fuzzy on the outskirts of the low-mass spiral galaxy NGC 247, 2023, MNRAS, 518, 3164}
\item G. Chen, H.-X. Zhang, X. Kong, et al. (including {\bf Y. Tang}), \href{https://ui.adsabs.harvard.edu/abs/2022ApJ...934L..35C/abstract}{Discovery of a Bimodal Environmental Distribution of Compact Ellipticals in the Local Universe, 2022, ApJL, 934, L35}
\item Y. Yao, H. Liu, X. Kong, et al. (including {\bf Y. Tang}), \href{https://ui.adsabs.harvard.edu/abs/2022ApJ...926...57Y/abstract}{Subgalactic Scaling Relations with $T_e$-based Metallicities of Low-metallicity Regions in Galaxies: Metal-poor Gas Inflow May Have Important Effects?, 2022, ApJ, 926, 57}
\item C. J. Agostino, S. Salim, S. M. Faber, et al. (including {\bf Y. Tang}), \href{https://ui.adsabs.harvard.edu/abs/2021ApJ...922..156A/abstract}{Physical Drivers of Emission-line Diversity of SDSS Seyfert 2s and LINERs after Removal of Contributions from Star Formation, 2021, ApJ, 922, 156}
\end{etaremune}

%-------------------------------------------
\end{document}
